% =============================================================================
% Discussion + Acknowledgments + Author Contributions + Conflicts of Interest +
% Funding + AI Use Disclosure + Data Availability Statement + Ethics Statement.
% Compiles via the canonical wrapper; see BUILD.md.
% =============================================================================


\section{Discussion}
\label{sec:discussion}

\subsection{Principal findings}
\label{discussion:findings}

No conventional evaluation approach exposes the four findings reported in \S\ref{sec:results} together: Simpson's paradox is invisible in aggregate-AUC reporting; direction asymmetry is invisible in unidirectional external validation; feature-transfer structure is invisible in discrimination-only evaluation; and algorithmic conditioning of the intraoperative advantage is invisible in aggregate reporting. Each requires a corresponding methodological commitment: \textbf{(1)} stratified within-cohort analysis across clinically meaningful risk strata; \textbf{(2)} bidirectional testing with multi-dimensional matched case-mix sensitivity; and \textbf{(3)} case-level paired bootstrap inference with rank-invariant metrics reported separately from probability-scale metrics. We operationalize these TRIPOD+AI-aligned commitments and argue for routine cross-population adoption. The sections that follow position each finding against the literature (\S\ref{discussion:literature}), examine the underlying mechanisms (\S\ref{discussion:mechanisms}), derive concrete methodological changes (\S\ref{discussion:deployment}), and consider the constraints that limit our conclusions (\S\ref{discussion:limitations}).

\subsection{Comparison with existing literature}
\label{discussion:literature}

The Simpson's paradox finding is, to our knowledge, the first systematic empirical demonstration in clinical prediction model evaluation. Although the paradox has been well-described in epidemiology since 1951~\cite{Simpson1951}, its implications for aggregate discrimination metrics have received limited attention~\cite{VanCalster2019}; the gap we report (overall AUC 0.756 masking near-random within-stratum AUCs 0.584--0.597) is the largest reported for a deployed clinical prediction model. The intraoperative-feature advantage extends the 5--10\% internal gain from vital sign features~\cite{Xue2021, Fritz2019} from internal to cross-population validation, preserved in 3 of 4 within-direction comparisons.

The direction-asymmetry finding extends a smaller cross-population literature: Futoma et al.~\cite{Futoma2020} and Finlayson et al.~\cite{Finlayson2021} characterized the ``myth of generalisability'' and distributional shift as core deployment risks but did not quantify direction-of-transfer asymmetry or provide a case-level paired bootstrap framework. The SHAP transferability finding (Spearman $\rho = 0.93$--$0.97$) is the first quantitative assessment of SHAP stability between geographically separated academic medical centers; prior SHAP-robustness work focused on within-dataset stability~\cite{Lundberg2017}, and our cross-population result suggests that the structure of what models learn is more robust than prior literature has demonstrated.

\subsection{Mechanisms}
\label{discussion:mechanisms}

The four-dimension matched-sensitivity analysis clarifies the mechanism underlying the observed direction asymmetry: ASA stratum matching amplifies the asymmetry to 128--203\% of baseline; Elixhauser comorbidity matching preserves or amplifies it (113--126\%); emergency-case proportion attenuates it to 70--86\% but does not eliminate it; and temporal period was untestable due to INSPIRE timestamp anonymization. The residual asymmetry surviving the three testable matches --- approximately 70\% of baseline --- is attributable to training-population characteristics beyond measured case-mix: healthcare-system-specific feature--outcome relationships, within-stratum severity coding that ASA and Elixhauser do not capture, and unmeasured temporal effects. At matched ASA stratum, INSPIRE deaths spread across the acuity spectrum (47.2\% / 52.8\% in ASA 1--2 / $\geq 3$) while MOVER deaths concentrate in high-acuity patients (1.1\% / 98.9\%) --- a 50.6-percentage-point split-difference consistent with cross-institutional feature--outcome relationships differing beyond what stratum matching equalizes. ASA physical status itself shows substantial inter-rater variability across institutions and clinicians~\cite{Sankar2014}, so a ``severity-3'' patient under INSPIRE coding may not be clinically equivalent to a ``severity-3'' patient under MOVER coding. Cross-institutional surveys of Asian-ICU end-of-life decision-making have documented systematic differences between Korean and Western practice patterns~\cite{Phua2015,Park2018}, consistent with the healthcare-system-specific mechanism interpretation above. Supplementary \S\ref{supp:s7-mechanism-evidence} reports stratified mortality distributions and the supporting cross-institutional literature.

The SHAP-plus-calibration finding provides a mechanistic lens on what transfers. Feature importance rankings are preserved almost perfectly ($\rho \geq 0.93$): a Seoul-trained model evaluated on Irvine data ranks feature importance (ASA, BMI, HR dynamics) in the same order as on its training cohort, yet absolute risk probabilities are systematically biased and require recalibration. Calibration fails uniformly before Platt scaling (O:E 0.02 to 2.92; slopes 0.41 to 1.29). This dissociation --- discrimination is preserved as a ranking, but calibration is distorted on the probability scale --- aligns with the well-documented pattern that discrimination transfers more readily than calibration~\cite{VanCalster2019}, supporting the interpretation that cross-population models preserve relational structure (which features matter, in what order) but not absolute mappings (what probability a given feature vector implies).

\subsection{Clinical implications}
\label{discussion:deployment}

These findings motivate three concrete changes to current practice. First, stratified external validation within clinically meaningful risk groups (e.g., ASA 1--2 vs $\geq 3$; elective vs emergency) should become standard for TRIPOD+AI-compliant reporting; aggregate AUC should accompany, not replace, within-stratum AUCs and paradox gaps (Supplementary Table~\ref{tab:paradox_summary}). Second, cross-population claims should include bidirectional transfer testing with multi-dimensional matched-case-mix sensitivity; $A \to B$ without reciprocal $B \to A$ cannot detect direction asymmetry, and matching is required to distinguish genuine failure from case-mix mismatch. Third, cross-population deployment requires mandatory probability recalibration (Platt or isotonic) on a representative sample of the deployment population. Recent external-validation sample-size guidance suggests that several thousand cases are typically needed to estimate calibration slope and O:E precisely, depending on the anticipated event proportion and degree of miscalibration~\cite{Riley2021Ext}. Discrimination-focused reporting alone is insufficient; published AUCs should accompany calibration assessments and recalibration protocols.

These methodological commitments are particularly consequential for cross-population deployment where training-population case-mix differs from deployment-population case-mix in ways conventional acuity classification does not capture. A salient instance is the deployment of models developed in high-resource academic medical centers to low- and middle-income-country (LMIC) tertiary settings. LMIC surgical populations are not uniformly lower-acuity than high-resource populations: depending on condition and local care-chain, patients may present at more advanced disease stages from delayed care-seeking, with longer-undertreated comorbidities and with selection-into-surgery patterns shaped by limited non-operative alternatives --- differences that ASA, comorbidity scores, and emergency-case proportion do not capture. Our data do not directly evaluate LMIC deployment, but the asymmetric transfer and residual matched-asymmetry patterns we report are precisely those such deployments would face; unidirectional validation studies would not surface this class of risk.

\subsection{Limitations}
\label{discussion:limitations}

Five principal limitations apply. First, only two institutions (Seoul National and UC Irvine) were included, characterizing cross-institutional rather than cross-national transfer. Second, temporal period matching was infeasible (INSPIRE distributes timestamps as privacy-preserving relative-time offsets); unmeasured temporal effects may contribute to the residual asymmetry. Third, neither race/ethnicity stratification (omitted from the public MOVER release; no Korean comparator) nor within-stratum severity matching beyond Elixhauser comorbidity could be uniformly operationalized; race/ethnicity is therefore not analyzed. This prevents detection of subgroup discrimination failures along race or ethnicity --- a well-documented fairness failure mode in cross-population deployment~\cite{Chen2020,Pfohl2021}. Fourth, procedure-type matching was not attempted because INSPIRE's KCD-10 and MOVER's CPT/EPIC code taxonomies do not align without a purpose-built crosswalk. Fifth, the 60-minute post-induction window was chosen by clinical convention; alternative lengths were not systematically examined. These limitations constrain the empirical reach of our perioperative findings; the methodological framework, however, applies independently of the specific application. Even with these constraints, the three methodological commitments are properties of the evaluation framework itself, and broader adoption would improve deployment safety across diverse populations and healthcare systems.

\section*{Acknowledgments}
The author acknowledges the Seoul National University Hospital INSPIRE team for stewardship of the INSPIRE dataset~\cite{Lim2024} and the UC Irvine perioperative team for stewardship of the MOVER dataset~\cite{Samad2023}.

\section*{Author Contributions}
Sole author. The author was responsible for all aspects of the work, including study conception and design, data acquisition under credentialed-access agreements, software implementation, formal analysis, interpretation, manuscript drafting, and final approval. The CRediT taxonomy roles --- conceptualization, methodology, software, validation, formal analysis, investigation, resources, data curation, writing (original draft), writing (review and editing), visualization, supervision, project administration, and funding acquisition --- all apply to the sole author.

\section*{Conflicts of Interest}
The author declares no competing interests.

\section*{Funding}
This work received no external funding.

\section*{AI Use Disclosure}
Anthropic's Claude (Claude Sonnet 4 and Claude Opus 4 series, accessed via Claude Code and the Claude.ai web interface, 2025--2026) was used during preparation of this manuscript for two purposes: (i) generation and refinement of analysis and figure-generation code in Python, which was reviewed and validated by the author before use; and (ii) editorial assistance with prose revision, copyediting, and structural organization of author-written drafts. All scientific claims, statistical analyses, study design choices, and interpretations are the author's; every quantitative result reported was independently validated against source data and primary literature by the author. No AI system is listed as an author, consistent with COPE and ICMJE guidance that AI tools cannot meet authorship criteria.

\section*{Data Availability Statement}
INSPIRE is publicly available via PhysioNet under credentialed access~\cite{INSPIRE_dataset}; MOVER is publicly available via the UC Irvine Machine Learning Repository under credentialed access~\cite{MOVER_dataset}. Both datasets were accessed under their respective data-use agreements; no patient-level data are redistributed in this paper. Analysis code, figure-generation scripts, trained model objects, and aggregate analysis outputs are archived at Zenodo (DOI assigned at submission via GitHub--Zenodo integration; see Methods \S\ref{methods:reproducibility}). Per-case prediction outputs are not redistributed (data-use agreement constraint); credentialed researchers can reproduce all analyses by applying the released models to source data after re-credentialing on PhysioNet and the UCI Machine Learning Repository.

\section*{Ethics Statement}
This study used de-identified data released for research use under the data-use agreements of INSPIRE (PhysioNet) and MOVER (UC Irvine Machine Learning Repository). No additional Institutional Review Board review was sought, as the work involves only secondary analysis of publicly released, de-identified datasets accessed under credentialed-use agreements. The author followed the data-use terms of both repositories.

% =============================================================================
% END OF DISCUSSION
% =============================================================================
